% ============================================================
\chapter{ビューの形：絶対・相対・順位・ファクター}\label{ch:shape}

同じ情報でも、ビューの形（行列 $P$）の選び方で、BL の結果は変わる。この章では、ビューの形の選択肢と、それぞれが何を変えるのかを整理する。とくに、LLM の予想のように「全体に楽観的すぎる」「全体に悲観的すぎる」といった癖がある予想では、相対ビューにするだけで、その癖を取り除ける。

\section{絶対ビュー}

1 つの資産の期待リターンを直接言うビューである。$p_k$ は、その資産の位置だけが 1 で、ほかは 0 の行ベクトルになる。対象論文と本プロジェクトは、全銘柄に絶対ビューを置き、$P=I$ としている。

絶対ビューには、銘柄ごとの差の情報と、全体の水準の情報が混ざっている。例えば、全銘柄の予想が一律に 0.1\% 高ければ、それは「市場全体が上がる」というビューである。BL はそれを、すべての銘柄の期待リターンを上げる方向に使い、制約のない最適化なら、市場ポートフォリオを買い増す（合計が 100\% を超える）。全額投資の制約があると、その情報は使い道がなく、ベータの高い銘柄に寄せるなど、意図しない形で比率に出る。

\section{相対ビュー}

2 つ以上の資産の差を言うビューである。$p_k$ の和は 0 で、上回る側に正、下回る側に負の重みを置く。相対ビューは、全体の水準についての情報を含まない。運用者の多くは、「A は B より良い」という形で考えることが多いので、相対ビューのほうが自然だと Idzorek は述べている。

重みの付け方は、\chapref{ch:idzorek} で見たとおり、時価総額の重み（Idzorek）と等しい重み（Satchell and Scowcroft）がある。時価総額の重みのほうが、小さい資産を不必要に大きく動かさずに済む。

\subsection*{LLM の予想を相対ビューにする}

対象論文の第 2 版は、LLM ごとに「投資スタイル」があり、それが成績を決めると結論している。論文の表 4 によると、回答の平均は、Gemma が $-0.38$\%、Qwen が $+0.10$\%、Llama が $+0.18$\%、GPT-4o-mini が $+0.13$\%（いずれも日次）で、Gemma は全体に悲観的だった。論文は、上昇相場では楽観的なモデルが勝ったと述べている。

これは、絶対ビューの「全体の水準」の部分が成績を決めた、ということである。全体の水準の予想は、市場全体の方向を当てる予想で、1 本の予想にすぎない（独立な予想の数が少ない）。それを数百のビューに分けて入れ、$\Omega$ を対角にすると、確信が数百倍に積み上がる（\secref{sec:diag-trap}）。

対策は、ビューから全体の水準を取り除くことである。

\begin{enumerate}
\item \textbf{平均を引く}：各時点で予想の平均（または時価総額加重の平均）を引き、$\tilde Q=Q-\bar Q\one$ を相対ビューとして使う。$P=I-\frac1N\one\one\T$（または時価総額の重みで引く形）になる。
\item \textbf{ペアやグループの差にする}：「上位 10\% の平均は下位 10\% の平均を上回る」のような、少数のロング・ショートのビューにまとめる。
\item \textbf{全体の水準は別の 1 つのビューにする}：市場全体の方向について言いたいことがあれば、市場ポートフォリオについての 1 つのビューとして、それに見合う不確かさで入れる。
\end{enumerate}

1 番目の $P=I-\frac1N\one\one\T$ は階数が $N-1$ で、$P\tau\Sigma P\T$ は逆行列を持たない。$\Omega$ が正定値なら、更新形の $(P\tau\Sigma P\T+\Omega)^{-1}$ は計算できる。

\subsection*{全体の水準が比率に効くのはいつか}

全体の水準の癖がいつ比率を動かすかは、$\Omega$ の形と、最適化の制約で決まる。

\begin{itemize}
\item \textbf{$\Omega\propto\Sigma$ で、全額投資の制約があるとき}：事後の期待リターンは $\Pi+s(Q-\Pi)$ で、$Q$ を一律に $c$ だけ上げると、すべての銘柄の期待リターンが $sc$ だけ一律に上がる。全額投資（比率の和が 1）なら、一律の上乗せは目的関数に定数を足すだけで、比率は変わらない。全体の水準の癖は、自動的に打ち消される（\chapref{ch:project} で確かめた）。
\item \textbf{$\Omega$ が対角のとき}：一律の上乗せ $c\one$ は、事後では $c\,\tau\Sigma(\tau\Sigma+\Omega)^{-1}\one$ になり、銘柄ごとに違う大きさで配られる。市場との共分散が大きい（ベータの高い）銘柄ほど大きく上がる。楽観的な予想は高ベータの銘柄への傾きに、悲観的な予想は低ベータの銘柄への傾きになる。上昇相場で楽観的な LLM が勝った、という対象論文の「スタイル」の結果は、この仕組みで説明できる。
\item \textbf{全額投資の制約がないとき}：全体の水準の上げ下げは、そのまま市場への投資額の増減になる。
\end{itemize}

相対ビューにする（平均を引く）のは、とくに $\Omega$ が対角の場合や、全額投資の制約がない場合に大切である。

\section{ポートフォリオへのビュー}

セクターや国、スタイルなど、資産のまとまりについてのビューである。$p_k$ はそのまとまりの比率（時価総額の重み）になる。「エネルギーのセクターは市場を 2\% 上回る」なら、エネルギー株の時価総額の比率から市場の比率を引いた行になる。

まとまりについてのビューは、個別の銘柄のビューより当たりやすいとは限らないが、個別の雑音が平均で消えるので、外れの分散は小さくなる。また、ビューの数が少ないので、$\Omega$ の相関の扱いも簡単になる。月次レポートのような、セクターや国の単位で書かれた情報は、この形にすると素直に入る（『総合報告書』の「月次レポートの読み取り」の章）。

\section{順位のビュー}

予想の値には自信がないが、順番には意味がある、という場合がある。Almgren と Chriss（2006）は、期待リターンの順序だけが分かっているときの最適なポートフォリオを考え、順序と矛盾しない期待リターンの集合の「中心」を使う方法を示した。Meucci（2010）は、順位のような「緩い」ビューを扱える方法として、Almgren and Chriss とエントロピー・プーリング（\chapref{ch:ep}）を挙げている。

実務でよく使われる簡単な方法は、順位を正規分布の点数に直し、Grinold の式で大きさを決めることである。
\begin{equation}
z_i=\Phi^{-1}\!\left(\frac{\mathrm{rank}_i-0.375}{N+0.25}\right),\qquad
\alpha_i=\IC\cdot\sigma_i\cdot z_i
\end{equation}
$\Phi^{-1}$ は標準正規分布の分位点（Blom の近似）。順位の当たり具合（順位相関の IC）だけで、期待リターンの大きさが決まる。LLM に数値ではなく順位や区分（強気・中立・弱気）を答えさせる場合にも使える。

\section{ファクターへのビュー}

資産の数が多いとき、すべての資産にビューを置くのは大変で、外れの相関の扱いも難しい。ファクター（市場、規模、バリュー、モメンタムなど）についてのビューにすれば、数が少なく、意味も分かりやすい。

\subsection*{ファクターモデルの BL}

資産のリターンがファクターモデル
\begin{equation}
r=Bf+\epsilon,\qquad f\sim\mathcal N(\mu_f,F),\qquad \epsilon\sim\mathcal N(0,D)
\end{equation}
に従うとする。$B$ は露出（$N\times L$）、$\mu_f$ はファクターのリスクプレミアムである。Kolm と Ritter（2017）は、裁定価格理論（APT）の文脈で、$\mu_f$ に事前分布 $\mathcal N(\xi,V)$ を置き、ファクターについてのビュー $P_f\mu_f=Q_f+\varepsilon_f$ で更新する形を詳しく扱っている。Spears ほか（2023）はこれを BL--APT と呼び、階層ベイズのモデルとして書き直した。

ファクターのプレミアムの事後分布は、ふつうの BL と同じ形である。
\begin{align}
\mu_f^{\mathrm{post}}&=\xi+VP_f\T(P_fVP_f\T+\Omega_f)^{-1}(Q_f-P_f\xi)\\
V^{\mathrm{post}}&=V-VP_f\T(P_fVP_f\T+\Omega_f)^{-1}P_fV
\end{align}
資産の期待リターンは $B\mu_f^{\mathrm{post}}$、予測の共分散は $B(F+V^{\mathrm{post}})B\T+D$ になる。Spears ほかは、この予測分布のもとでの最適な比率を、行列の公式を使って閉じた形で書いている。

ファクターへのビューの利点は次のとおりである。

\begin{itemize}
\item ビューの数が少ない（数個〜数十個）。外れの相関も、ファクターの共分散 $F$ を使えば自然に入る。
\item 1 つのビューが、露出 $B$ を通じて全資産に一貫した形で効く。
\item ファクターのプレミアムには長い歴史のデータがあるので、事前分布 $\xi$、$V$ と、ビューの当たり具合を測りやすい。
\end{itemize}

数値例は \chapref{ch:esg} で、セクター ETF と Fama--French のファクターを使って示す。

\subsection*{ほかの拡張}

Meucci（2009）は、資産のリターンではなく、一般的なリスクファクター（金利、ボラティリティなど、損益に非線形に効くもの）にビューを置き、ストレステストにも使える形に拡張した。Cheung（2013）の Augmented BL は、ファクターへのビューを中心にした BL で、ファクターのランキングに頼らずに、ファクターに基づくポートフォリオを作る方法を示している。

\section{外部の変数へのビュー}

ビューを置く対象は、投資する資産でなくてもよい。Meucci（2010）は、マクロ経済の変数のような「外部の要因」についてのビューを、相関を通じて資産のリターンに反映させる使い方を、BL では扱えず COP やエントロピー・プーリングなら扱える例として挙げている。実は、市場版の BL でも、外部の変数 $g$ を資産のリターン $X$ と一緒に並べた $(X,g)$ の正規分布を考えれば、$g$ へのビューで $X$ を条件付けられる。

\begin{equation}
\E[X\mid g=v]=\pi_X+\Sigma_{Xg}\,\Sigma_{gg}^{-1}(v-\pi_g)
\end{equation}
これは、$g$ への回帰の係数で $X$ を動かすことと同じである。\chapref{ch:politics} で、地政学リスク指数の変化を外部の変数にした例を示す。

\section{ビューの数と、ビューどうしの関係}

\begin{itemize}
\item \textbf{ビューの数}：$K\le N$ である必要はないが、$K$ が多いほど、外れの相関の扱いが重要になる。
\item \textbf{同じことを言うビュー}：例えば「A は 5\%」「B は 3\%」「A$-$B は 2\%」の 3 つは、3 つ目が前の 2 つから決まる。外れが独立だと扱うと、同じ情報を 2 回数えることになる。
\item \textbf{矛盾するビュー}：「A$-$B は 2\%」「A$-$B は $-1\%$」のように食い違うビューは、不確かさに応じた平均になる。食い違いが大きいなら、どちらかの不確かさを見直す（\chapref{ch:text} の融合の方法）。
\item \textbf{計算の安定性}：$P\Sigma P\T$ がほぼ特異（ビューのポートフォリオがほとんど同じ）だと、100\% に近い確信で計算が不安定になる。
\end{itemize}

\begin{tipbox}[コツ：ビューの形]
\begin{enumerate}
\item 全体の水準と銘柄間の差は、別々に考える。銘柄間の差は相対ビュー（平均を引く）、全体の水準は市場への 1 つのビューにする。
\item 相対ビューの重みは時価総額で付ける。
\item 予想の値に自信がないときは、順位を点数に直し、IC で大きさを決める。
\item 資産が多いときは、ファクターへのビューを検討する。数が少なく、相関も自然に入る。
\item 同じ情報を 2 回数えていないか、ビューを並べて確かめる。
\end{enumerate}
\end{tipbox}

\begin{keybox}[この章のまとめ]
\begin{itemize}
\item 絶対ビューは、銘柄間の差と全体の水準を混ぜて持つ。LLM の「スタイル」（全体の楽観・悲観）は、この全体の水準の部分である。
\item 相対ビュー（平均を引く、ペアにする）にすれば、全体の水準の癖を取り除ける。
\item ファクターへのビューは、数が少なく、相関が自然に入り、データで当たり具合を測りやすい。
\item 外部の変数へのビューも、資産と一緒の正規分布を考えれば、市場版の BL で扱える。
\end{itemize}
\end{keybox}
